\begin{recipe}
[% 
    preparationtime = {\unit[20]{min}},
    portion = {\portion{2}},
    source = {\protect\recipesource{https://www.instagram.com/p/DDh6R8us75L/}{Instagram: fit\_sophia\_laetitia}},
    calory = {588kcal (938kcal) | 23g (30g) Protein | 57g (134g) KH | 28g (29g) Fett},
    price = {2,90€ (3,50€)},
    %rating = {1},
]
{Schnelles Linsen Curry}

%\graph{% pictures
%    big=pic/schnellekueche/02_linsen-curry
%}

\introduction{%
Ein cremiges, würziges Linsen-Curry, das in unter 20 Minuten auf dem Tisch steht. Viel Protein, viel Geschmack und perfekt für stressige Tage.
}

\ingredients{
    2 & Schalotten\index{Gemüse!Schalotte}\\
    1 Stück & Ingwer (daumengroß)\index{Gemüse!Ingwer}\\
    3 Zehen & Knoblauch\index{Gemüse!Knoblauch}\\
    2 Dosen & Linsen (abgetropft ca. 480 g)\index{Hülsenfrüchte!Linsen}\\
    200 ml & Kokosmilch\index{Milchprodukte \& Alternativen!Kokosmilch}\\
    300 ml & Passierte Tomaten\index{Gemüse!Tomate}\\
    2 EL & Tomatenmark\\
    1 EL & Ghee\\
    1 TL & Kreuzkümmel\\
    1 TL & Paprikapulver (rosenscharf)\\
    1 TL & Curry\\
    1 TL & Kurkuma\\
    \nicefrac{1}{2} TL & Salz\\
    & Salz und Zitronensaft zum Abschmecken
}

\preparation{%
    \step%
    Schalotten, Knoblauch und Ingwer fein hacken.\\
    
    \step%
    Ghee in einer Pfanne erhitzen und alles 2--3 Minuten glasig anbraten.
    
    \step%
    Linsen, Kokosmilch, Tomaten, Tomatenmark und Gewürze zugeben und 8--10 Minuten einkochen lassen, bis die gewünschte Konsistenz erreicht ist.
    
    \step%
    Mit Salz und etwas frischem Zitronensaft abschmecken.
}

\suggestion[Servieren]{%
    Mit Chiliöl, etwas zusätzlicher Kokosmilch, frischem Koriander und frischer Chili servieren.
    Mit 100 g ungekochtem Reis pro Portion servieren. 
}

\end{recipe}